%% Les trois structures algorithmiques de base : linéaire, alternative, répétitive.
\documentclass[tikz,border=4pt]{standalone}
\usepackage{liaisons}
\begin{document}
\begin{tikzpicture}[x=1cm,y=1cm,
  trait/.style={draw=solideD, line width=1.1pt, fill=solideD!12, rounded corners=1pt,
                minimum width=1.5cm, minimum height=0.55cm, inner xsep=2pt, inner ysep=1pt,
                align=center, font=\scriptsize},
  los/.style={draw=solideB, line width=1.1pt, fill=white, line join=round},
  fl/.style={-{Stealth[length=4.5pt]}, line width=0.8pt},
  boucle/.style={-{Stealth[length=4.5pt]}, line width=0.9pt, solideF, rounded corners=3pt},
  titre/.style={font=\small\bfseries, anchor=south},
  bal/.style={font=\scriptsize, inner sep=1.5pt}]

%% ================== 1. structure linéaire =============================
\begin{scope}[shift={(0,0)}]
  \node[titre] at (0,1.15) {Linéaire};
  \node[trait] (l1) at (0,0.4)  {Traitement 1};
  \node[trait] (l2) at (0,-0.7) {Traitement 2};
  \node[trait] (l3) at (0,-1.8) {Traitement 3};
  \draw[fl] (0,0.95) -- (l1.north);
  \draw[fl] (l1.south) -- (l2.north);
  \draw[fl] (l2.south) -- (l3.north);
  \draw[fl] (l3.south) -- (0,-2.55);
\end{scope}

%% ================== 2. structure alternative ==========================
\begin{scope}[shift={(2.9,0)}]
  \node[titre] at (0,1.15) {Alternative};
  \draw[los] (0,-0.45) -- (0.85,0) -- (0,0.45) -- (-0.85,0) -- cycle;
  \node[font=\scriptsize] at (0,0) {Condition};
  \draw[fl] (0,0.95) -- (0,0.47);
  % branche « oui » vers le bas, branche « non » vers la droite
  \node[trait] (a1) at (-0.9,-1.35) {Traitement 1};
  \node[trait] (a2) at (0.9,-1.35)  {Traitement 2};
  \draw[fl] (0,-0.45) -- (0,-0.75) -- (-0.9,-0.75) -- (a1.north);
  \node[bal, text=solideE, anchor=north east] at (-0.04,-0.48) {oui};
  \draw[fl] (0.85,0) -- (0.9,0) -- (0.9,-0.75) -- (a2.north);
  \node[bal, text=solideA, anchor=south] at (0.92,-0.02) {non};
  \draw[line width=0.8pt] (a1.south) -- (-0.9,-2.15) -- (0.9,-2.15) -- (a2.south);
  \draw[fl] (0,-2.15) -- (0,-2.6);
\end{scope}

%% ================== 3. structure répétitive ===========================
\begin{scope}[shift={(6.0,0)}]
  \node[titre] at (1.0,1.15) {Répétitive};
  \draw[fl] (0,1.1) -- (0,0.64);
  \fill[solideF] (0,0.6) circle (1.6pt);
  \draw[fl] (0,0.6) -- (0,0.47);
  \draw[los] (0,-0.45) -- (0.85,0) -- (0,0.45) -- (-0.85,0) -- cycle;
  \node[font=\scriptsize] at (0,0) {Condition};
  \node[trait] (r1) at (2.05,0) {Traitement};
  \draw[fl, solideE] (0.85,0) -- (r1.west);
  \node[bal, text=solideE, anchor=south] at (1.06,0.02) {oui};
  \draw[boucle] (r1.north) -- (2.05,0.6) -- (0.06,0.6);
  \draw[fl] (0,-0.45) -- (0,-1.15);
  \node[bal, text=solideA, anchor=west] at (0.04,-0.8) {non};
\end{scope}
\end{tikzpicture}
\end{document}
